% P15.21: Feller VII.7 Theorem 2 and Durrett Examples 2.2.15--2.2.16.
\chapter{Weak Normalization with Infinite Means}
\label{sec:p15-infinite-mean-weak-normalization}

A positive gain can have infinite expectation while the total from many
independent trials concentrates around a deterministic scale. The scale
then need not be proportional to the number of trials. This chapter
gives Feller's necessary and sufficient condition, constructs its
normalizers, and retains the separate triangular-array calculation for
the St.~Petersburg game. The sources are Feller, Volume II, second
edition, VII.7, Theorem~2 and its continuation
\cite[pp.~236--237]{Feller1971ProbabilityII}, and Durrett's author
Version~5 dated 11 January 2019, Examples~2.2.15--2.2.16
\cite[pp.~65--66]{DurrettPTEAuthor2019}.

\section{Positive gains and exact truncation conventions}
\label{sec:p15-positive-gain-conventions}

\begin{hbdefinition}{def:p15-positive-gain-normalization}{The law, its tail and its truncated moments}
Let \(X_1,X_2,\ldots\) be independent identically distributed finite
real random variables on \((\Omega,\mathcal F,\mathbb P)\), with
\(X_1\geq0\) almost surely and \(\mathbb P(X_1>0)>0\).
Write \(X\) for a variable with this common law,
\(F(s)=\mathbb P(X\leq s)\), and \(S_n=\sum_{k=1}^nX_k\).
Feller's stated domain is the special case \(\mathbb P(X>0)=1\),
equivalently \(F(0)=0\). The results below also prove the indicated
extension allowing an atom at zero.

For \(s>0\), define
\[
 q(s)=\mathbb P(X>s),\quad
 \mu(s)=\mathbb E[X\mathbf1_{\{X\leq s\}}],\quad
 v(s)=\mathbb E[X^2\mathbf1_{\{X\leq s\}}],
\]
and
\[
 \mu_<(s)=\mathbb E[X\mathbf1_{\{X<s\}}],\qquad
 q_{\geq}(s)=\mathbb P(X\geq s).
\]
These moments are finite even when \(\mathbb EX=\infty\).
Choose any \(s_0>0\) with \(\mu(s_0)>0\). For \(s\geq s_0\), put
\[
 \delta(s)=\frac{s q(s)}{\mu(s)},\qquad
 \rho(s)=\frac{\mu(s)}{s q(s)}.
\]
If \(q(s)=0\), set \(\rho(s)=+\infty\) and \(\delta(s)=0\).
Thus \(\rho(s)\to\infty\) means exactly \(\delta(s)\to0\).
No ratio with zero truncated mean is used.
A \emph{relative weak normalizer} is a deterministic sequence
\(a_n>0\) such that \(S_n/a_n\to1\) in probability:
\(\mathbb P(|S_n/a_n-1|>\varepsilon)\to0\) for every
\(\varepsilon>0\).
\index{weak law!infinite mean}\index{normalization!relative weak}
\end{hbdefinition}

\begin{hblemma}{thm:p15-positive-truncation-tools}{proved=annex:p15.positive-truncation-tools}{Truncated moment identities and the two-error bound}
For every \(s>0\),
\[
 \mu(s)=\int_0^s q(t)\,dt-sq(s),\qquad
 v(s)=2\int_0^s t q(t)\,dt-s^2q(s),\qquad
 v(s)\leq s\mu(s).
\]
For \(a>0\), \(\varepsilon>0\) and \(n\geq1\),
\[
 \mathbb P(|S_n-n\mu(s)|>\varepsilon a)
 \leq \frac{n v(s)}{\varepsilon^2a^2}+n q(s).
\]
Moreover \(\mu(s)/s\to0\). If \(\delta(s)\to0\), then
\[
 \frac{v(s)}{s\mu(s)}\longrightarrow0.
\]
The deletion truncation retains the atom at \(s\); the error event
uses the strict tail \(X>s\).
\end{hblemma}

\begin{hblemma}{thm:p15-relative-normalizer-properties}{proved=annex:p15.relative-normalizer-properties}{Growth, adjacent ratios and uniqueness of scale}
Under the positive-gain assumptions, every relative weak normalizer
satisfies \(a_n\geq c n\) for all sufficiently large \(n\), for some
\(c>0\), and
\[
 a_n\longrightarrow\infty,\qquad
 \frac{a_{n+1}}{a_n}\longrightarrow1.
\]
No monotonicity of \(a_n\) is assumed. If \(a_n\) and \(b_n\) both
normalize \(S_n\) relatively to one, then \(a_n/b_n\to1\).
Conversely, replacing a relative weak normalizer by positive
\(b_n\) with \(b_n/a_n\to1\) preserves the relative weak limit.
\end{hblemma}

\section{The normalizer construction and Feller's criterion}
\label{sec:p15-feller-positive-criterion}

\begin{hblemma}{thm:p15-positive-normalizer-construction}{proved=annex:p15.positive-normalizer-construction}{An exact anchored crossing}
Choose \(s_0>0\) with \(\mu(s_0)>0\), and choose an integer \(N_0\)
so that \(n\mu(s_0)>s_0\) for \(n\geq N_0\).
For those integers define
\[
 s_n=\inf\{s\geq s_0:n\mu(s)/s\leq1\}.
\]
The set is nonempty, its infimum is finite and exceeds \(s_0\), and
\[
 n\mu(s_n)=s_n,\qquad
 s_n\geq n\mu(s_0)\longrightarrow\infty.
\]
The sequence \(s_n\), \(n\geq N_0\), is nondecreasing.
These existence statements require no finite expectation and no
assumption on \(\delta\). Arbitrary positive values may be assigned
at the finitely many indices below \(N_0\).
\end{hblemma}

\begin{hbdefinition}{def:p15-positive-slow-variation}{Slow variation of the truncated mean}
An eventually positive function \(L:(0,\infty)\to[0,\infty)\)
is \emph{slowly varying at infinity} if
\(L(cs)/L(s)\to1\) as \(s\to\infty\) for every fixed \(c>0\).
Here this definition is applied to the nondecreasing truncated
mean \(\mu\), with its eventual positive domain understood.
\index{slow variation!truncated mean}
\end{hbdefinition}

\begin{hbtheorem}{thm:p15-positive-tail-equivalences}{proved=annex:p15.positive-tail-equivalences}{Slow variation and either endpoint convention}
Under the positive-gain assumptions, the following are equivalent:
\[
 \delta(s)\to0;\qquad
 \mu\text{ is slowly varying};\qquad
 \frac{s q_{\geq}(s)}{\mu_<(s)}\to0.
\]
The last ratio is taken only on the eventual domain
\(\mu_<(s)>0\). Under these conditions,
\[
 \frac{\mu_<(s)}{\mu(s)}\to1,\qquad
 \frac{s\,\mathbb P(X=s)}{\mu(s)}\to0.
\]
In the infinite-mean case this is exactly the truncated-first-moment
specialization of the relation noted in Feller's footnote to VII.7,
Theorem~2; the needed equivalence is proved here directly
\cite[pp.~236, 282--284]{Feller1971ProbabilityII}.
The finite-mean case is included as well.
\end{hbtheorem}

\begin{hbtheorem}{thm:p15-feller-positive-weak-normalization}{proved=annex:p15.feller-positive-weak-normalization}{Feller VII.7, Theorem 2: both directions}
For independent identically distributed nonnegative gains with
\(\mathbb P(X>0)>0\), there exist positive deterministic \(a_n\)
with \(S_n/a_n\to1\) in probability if and only if
\[
 \rho(s)\longrightarrow\infty\quad(s\to\infty).
\]
Allowing arbitrary nonzero real constants does not enlarge the
conclusion: any sequence giving this limit is eventually positive.
When the condition holds, the anchored construction gives
\[
 a_n=n\mu(s_n)=s_n\qquad(n\geq N_0)
\]
and, for every \(\varepsilon>0\),
\[
 \mathbb P(|S_n/s_n-1|>\varepsilon)
 \leq \frac{v(s_n)}{\varepsilon^2s_n\mu(s_n)}
                              +\delta(s_n)\longrightarrow0.
\]
Feller's sufficiency argument uses the truncated second moment.
His necessity argument uses positivity, the largest positive
summand, and truncation at \(2a_n\). Both arguments, including the
passage from the sequence \(2a_n\) to every large real threshold,
are supplied in the annex.
\end{hbtheorem}

\begin{hbcorollary}{thm:p15-positive-normalized-centers}{proved=annex:p15.positive-normalized-centers}{All relative normalizers and their truncated centers}
Suppose the equivalent conditions of the preceding theorem hold,
and let \(s_n\) be its constructed normalizer. A positive sequence
\(a_n\) is a relative weak normalizer if and only if \(a_n/s_n\to1\).
For every such \(a_n\),
\[
 \frac{n\mu(a_n)}{a_n}\to1,\quad
 \frac{n\mu_<(a_n)}{a_n}\to1,\quad
 n q(a_n)\to0,\quad n q_{\geq}(a_n)\to0.
\]
Thus the closed and strict truncated centers have the same leading
scale, although only the closed mean is used in the exact
right-continuous crossing construction.
If \(0<\mathbb EX=m<\infty\), the canonical normalizer satisfies
\(s_n/n\to m\), and \(nm\) is a relative weak normalizer.
If \(\mathbb EX=\infty\), then \(s_n/n\to\infty\).
\end{hbcorollary}

\section{Worked positive laws and required hypotheses}
\label{sec:p15-positive-normalization-examples}

\begin{hbtheorem}{thm:p15-positive-normalization-boundaries}{proved=annex:p15.positive-normalization-boundaries}{Finite means, a critical tail and failures}
\begin{enumerate}
\item If \(X=c>0\) almost surely and \(s_0=c\), then \(s_n=nc\)
for \(n\geq2\). If \(X\) equals \(1\) or \(3\), each with
probability \(1/2\), the choice \(s_0=3\) gives \(s_n=2n\)
for \(n\geq2\).
\item For the probability density
\(f(x)=x^{-2}\mathbf1_{\{x\geq1\}}\),
\[
 q(s)=s^{-1},\quad \mu(s)=\log s,\quad
 v(s)=s-1\qquad(s\geq1).
\]
Its mean is infinite and its relative weak normalizers satisfy
\(s_n\sim n\log n\), where these logarithms are natural.
\item For \(0<\alpha<1\), the probability density
\(f_\alpha(x)=\alpha x^{-\alpha-1}\mathbf1_{\{x\geq1\}}\)
has infinite mean and
\[
 \mu(s)=\frac{\alpha}{1-\alpha}(s^{1-\alpha}-1),\qquad
 \rho(s)\longrightarrow\frac{\alpha}{1-\alpha}.
\]
It therefore has no relative weak normalizer.
\item If all \(X_k\) are the same variable taking values \(1,3\)
with equal probabilities, their marginal law meets the tail
criterion, but no relative weak normalizer exists.
The identically zero law likewise has no normalizer giving limit one;
it is excluded by \(\mathbb P(X>0)>0\).
\end{enumerate}
Here \(u_n\sim v_n\) means \(u_n/v_n\to1\) for positive quantities.
\end{hbtheorem}

\section{The Cauchy boundary and St. Petersburg}
\label{sec:p15-cauchy-and-petersburg}

Each of the following source examples specifies its own common law
and reuses \(S_n=\sum_{k=1}^nX_k\) for that law. The Cauchy variables
are signed; the positive-gain hypotheses apply again to
St.~Petersburg variables.

\begin{hbdefinition}{def:p15-cauchy-petersburg-laws}{The two source laws}
For \(a>0\), the centered Cauchy law with scale \(a\) has density
\[
 f_a(x)=\frac{a}{\pi(a^2+x^2)}\qquad(x\in\mathbb R).
\]
The standard Cauchy law has \(a=1\).
The St.~Petersburg gain law has masses
\[
 \mathbb P(X=2^j)=2^{-j},\qquad j=1,2,\ldots.
\]
It is realized by independent fair coin tosses, paying \(2^j\) if the
first head is on toss \(j\). The event of no head has probability zero.
Repeated plays use independent copies of the gain.
In the St.~Petersburg statements below, \(\log_2\) always denotes
base-two logarithm.
\index{Cauchy distribution}\index{St. Petersburg paradox}
\end{hbdefinition}

\begin{hbtheorem}{thm:p15-cauchy-weak-law-failure}{proved=annex:p15.cauchy-weak-law-failure}{Durrett Example 2.2.15 and its distributional assertion}
The Cauchy densities integrate to one. For a standard Cauchy \(X\),
\[
 \mathbb P(|X|>x)\sim\frac{2}{\pi x}\quad(x\to\infty),
 \qquad \mathbb E|X|=\infty.
\]
For independent standard Cauchy \(X_k\), no real deterministic
sequence \(c_n\) satisfies \(S_n/n-c_n\to0\) in probability.
In fact \(S_n/n\) has exactly the standard Cauchy law for every
\(n\geq1\). More generally \(f_a*f_b=f_{a+b}\) almost everywhere
as densities for \(a,b>0\).
The tail argument is Durrett's proof in Example~2.2.15.
The annex additionally proves its later distributional assertion by
direct real convolution, without a characteristic-function premise.
\end{hbtheorem}

\begin{hbtheorem}{thm:p15-petersburg-array-calculation}{proved=annex:p15.petersburg-array-calculation}{Durrett Example 2.2.16: the complete truncation route}
Let \(X_k\) be independent St.~Petersburg gains.
Their common expectation is infinite. Choose real \(K(n)\to\infty\)
so that \(m(n)=\log_2 n+K(n)\) is an integer, and put
\[
 b_n=2^{m(n)}=n2^{K(n)},\qquad
 \overline X_{n,k}=X_k\mathbf1_{\{X_k\leq b_n\}},\qquad
 a_n=nm(n).
\]
For all sufficiently large \(n\), \(m(n)\geq1\), and
\[
 \begin{aligned}
 n\mathbb P(X>b_n)&=2^{-K(n)},&
 n\mathbb P(X\geq b_n)&=2^{1-K(n)},\\
 \mathbb E\overline X_{n,k}&=m(n),&
 \mathbb E\overline X_{n,k}^{\,2}&=2^{m(n)+1}-2\leq2b_n.
 \end{aligned}
\]
Consequently \((S_n-a_n)/b_n\to0\) in probability.
Whenever \(K(n)/\log_2n\to0\), the deterministic center satisfies
\(a_n/(n\log_2n)\to1\).
If, in addition, \(K(n)\leq\log_2\log_2n\) eventually, then
\[
 \frac{S_n}{n\log_2n}\longrightarrow1
                                      \quad\hbox{in probability}.
\]
All these conditions are simultaneously met, for \(n\geq2\), by
\[
 m(n)=\left\lfloor\log_2(n\log_2n)\right\rfloor,\qquad
 K(n)=m(n)-\log_2n.
\]
The floor \(\lfloor u\rfloor\) is the greatest integer at most \(u\).
Initial finitely many choices do not affect any assertion.
\end{hbtheorem}

\begin{hbtheorem}{thm:p15-petersburg-normalizers-and-prices}{proved=annex:p15.petersburg-normalizers-and-prices}{Feller's route, finite calculations and the meaning of the scale}
For St.~Petersburg gains and \(s\geq2\), put
\(\ell(s)=\lfloor\log_2s\rfloor\). Then
\[
 \mu(s)=\ell(s),\quad
 v(s)=2^{\ell(s)+1}-2,\quad
 q(s)=2^{-\ell(s)},\quad
 \delta(s)\leq\frac{2}{\ell(s)}\longrightarrow0.
\]
Thus Feller's separate construction with \(s_0=2\) gives exact
normalizers \(s_n=n\lfloor\log_2s_n\rfloor\), for all sufficiently
large \(n\), and \(s_n/(n\log_2n)\to1\).
For the explicit array choice in the preceding theorem, \(n=1024\)
gives \(m(n)=13\), \(b_n=8192\), \(a_n=13312\), and
\[
 n\mathbb P(X>b_n)=\tfrac18,\qquad
 \frac{n\mathbb E\overline X_{n,k}^{\,2}}{b_n^2}
                   =\frac14-\frac1{32768}.
\]
The relative weak limit assigns the leading aggregate scale
\(n\log_2n\), or \(\log_2n\) per play. At \(n=1024\) this
per-play scale is \(10\). The scale equals \(40\) when
\(n=2^{40}=1\,099\,511\,627\,776\).
These are asymptotic scale calculations; they neither equal the
infinite expectation nor specify a finite-sample break-even probability.
For every fixed fee \(c>0\), however,
\(\mathbb P(S_n>cn)\to1\).

No deterministic \(c_n\) makes \(S_n/n-c_n\to0\) in probability:
at \(x=2^j\), \(x\mathbb P(X>x)=1\).
Changing the denominator is essential.
\end{hbtheorem}

\begin{proofboundary}{def:p15-positive-normalization-boundary}{Exact continuation}
This chapter completes Feller VII.7, Theorem~2, its required
slow-variation equivalence, and Durrett Examples~2.2.15--2.2.16.
The general regular-variation theorem of Feller VIII.9 has wider
moment orders and endpoint cases than the specialization proved here.
Its full presentation, the alternate argument in XVII.2a, and the
separately identified Poisson/Bessel example remain separate source
components. The complete numbered weak-convergence presentation is
\cref{sec:p15-weak-convergence-of-laws}, including its quantile,
portmanteau, selection and metric arguments.
The distinct characteristic-function proof of the Cauchy transform
and average law is supplied in
\cref{thm:p15-characteristic-inverse-examples}.
The direct convolution proof here retains its separate route.
The complete remaining characteristic-function source chain is
\cref{sec:p15-characteristic-function-criteria-and-moments}.
The central-limit constructions and exercises are supplied in
\cref{sec:p15-central-limit-theorems}. The complete normal-domain
criterion is \cref{thm:p15-nda-characterization}, with its original
truncated-mean cutoff presentation in \cref{thm:p15-nda-gk-cutoff-route}.
The later full stable-law presentations retain their distinct source
obligations.
\end{proofboundary}
